\begin{proof}
Choose a very ample integral toric divisor \(A\) with full-dimensional
polytope \(Q\). The convex dictionary is additive under addition of
toric currents, hence
\[
               \Delta(T+\varepsilon\omega_A)
                         =K+\varepsilon Q.
\tag{69.28}\label{eq:verified-minkowski-perturbation}
\]
The complete character-by-character calculation is given in the proof of
\Cref{thm:ok-toric-zero-mass}; we apply it to the present notation.  More
precisely, that proof converts the book's squared-norm weight to
Guenancia's norm weight, identifies the interiors of the local Newton set
and its closure, and checks that the strict condition automatically makes
the local Laurent exponent holomorphic.  After those checks, Guenancia's
criterion \cite[Theorem~1.13]{Guenancia12} says on a maximal smooth chart
\(U_\sigma\) that the character \(\chi^u\) survives exactly when
\[
 u+\kappa_\sigma\in
 \operatorname{int}\bigl(m(K+\varepsilon Q)+\sigma^\vee\bigr),
\]
where \(\kappa_\sigma\) is the fixed all-ones vector in the chart lattice.
The proof cited above checks the Cartier-data signs, proves the global gluing
identity
\[
 K+\varepsilon Q
 =\bigcap_{\sigma\in\Sigma(n)}
   (K+\varepsilon Q+\sigma^\vee),
\]
and squeezes the strict conditions with shifts \(m^{-1}\kappa_\sigma\) between
inner and outer convex bodies.  The lattice Riemann-sum theorem then proves
that the following inner limit exists and that
\[
                         S_A(\varepsilon)
                         =\vol_n(K+\varepsilon Q),
\tag{69.29}\label{eq:verified-toric-rr}
\]
where \(\vol_n\) is normalized by the character lattice.  Thus in this toric
case the normalized section count has an actual limit.  The displayed equality
holds only after that limit, not as a literal finite-level section-count
identity, and no factor \(n!\) occurs.

Put \(d=\dim K\) and let \(A_K\) be its affine span. After translation,
decompose
\[
                  M_{\mathbf R}=V_K\oplus V_K^\perp,
       \qquad \dim V_K=d.
\]
Because \(K\) has nonempty relative interior in \(A_K\), it contains a
\(d\)-dimensional ball of radius \(a>0\) in \(V_K\). Since \(Q\) has
nonempty interior, it contains an \(n\)-dimensional ball of radius
\(b>0\). Therefore \(K+\varepsilon Q\) contains a product-type body
with \(d\) directions of fixed size and \(n-d\) directions of size
\(\varepsilon\). Hence
\[
                \vol_n(K+\varepsilon Q)
                       \geq c\varepsilon^{\,n-d}.
\tag{69.30}\label{eq:verified-volume-lower}
\]

Conversely, \(K\) is contained in a bounded \(d\)-dimensional box in
\(A_K\), and \(Q\) is contained in a bounded \(n\)-dimensional box.
Their Minkowski sum is contained in a box with \(d\) sides bounded
independently of \(\varepsilon\) and \(n-d\) sides bounded by
\(C\varepsilon\). Thus
\[
                \vol_n(K+\varepsilon Q)
                       \leq C'\varepsilon^{\,n-d}.
\tag{69.31}\label{eq:verified-volume-upper}
\]
Equations \eqref{eq:verified-cao-characterization}--
\eqref{eq:verified-volume-upper} show that the exponent in Cao's
criterion is \(n-d\). Therefore
\(\nu_{\mathrm{Cao}}(T)=d\).
\end{proof}
